%! Simplifying Rational Expressions — Unit 1, Lesson 1.7 — Group Activity (Answer Key)
\documentclass[10pt]{article}
\usepackage{atda-article}
\usepackage{atda-key}   % pulls in -boxes; do NOT also load -boxes
\newcommand{\TallMath}[1]{$\displaystyle #1\rule[-1.4em]{0pt}{3.2em}$}
\setlength{\workrowsep}{4pt}

\begin{document}

\pageheader{Unit 1, Lesson 1.7}{Group Activity --- Answer Key}
% NAMESTRIP: no \namedateperiod here — mirrors the blank. Teacher-only prose goes
% in the lesson plan as \begin{teachernote}[Group Activity], never in a key.

\vspace{0.15in}

\begin{scenariobox}[The Scenario]{royal}
\small
The drama club is building a rectangular backdrop for the spring musical. The shop teacher measures
the finished panel and reports its area as
\[ A(x) = x^2+3x-28 \quad \text{square feet,} \]
and the frame drawing says the backdrop's \textbf{width} is $(x-4)$ feet. Nobody wrote down the
length, and the club has to order edge trim. Every question today is the same question:
\textbf{factor, then divide out what the top and bottom share.}

Work with your group. \textbf{Everyone starts at Tier R.} When your whole group can do Tier R
without help, move on.
\end{scenariobox}

\vspace{0.10in}

\boxguard[16]
\begin{tcolorbox}[colback=white, colframe=black!40, title=\textbf{Tier R --- Remediate},
    fonttitle=\bfseries, arc=2mm, left=3mm, right=3mm, top=2mm, bottom=2mm, breakable]
\small
Items 1--2: state every excluded value. Items 3--5: simplify completely \emph{and} state every
excluded value, read from the original denominator.

\begin{enumerate}[leftmargin=1.6em, itemsep=4pt, topsep=4pt]
  \item \TallMath{\frac{7}{x-5}}
\begin{work}
  x-5 &= 0 \quad\Rightarrow\quad x = 5 \quad \text{(excluded)}
\end{work}

  \item \TallMath{\frac{x+1}{x^2-9}}
\begin{work}
  (x+3)(x-3) &= 0 \\
  x &= -3 \quad \text{or} \quad x = 3 \quad \text{(both excluded)}
\end{work}

  \item \TallMath{\frac{6x^3}{8x}}
\begin{work}
  \frac{6x^3}{8x} &= \frac{2x \cdot 3x^2}{2x \cdot 4} = \frac{3x^2}{4} \qquad (x \ne 0)
\end{work}

  \item \TallMath{\frac{x^2-4}{x^2+5x+6}}
\begin{work}
  \frac{x^2-4}{x^2+5x+6} &= \frac{(x+2)(x-2)}{(x+2)(x+3)} \\
      &= \frac{x-2}{x+3} \qquad (x \ne -2, \; x \ne -3)
\end{work}

  \item \TallMath{\frac{3x+12}{x^2-16}}
\begin{work}
  \frac{3x+12}{x^2-16} &= \frac{3(x+4)}{(x+4)(x-4)} \\
      &= \frac{3}{x-4} \qquad (x \ne -4, \; x \ne 4)
\end{work}
\end{enumerate}
\end{tcolorbox}

\vspace{0.10in}

\boxguard[16]
\begin{tcolorbox}[colback=white, colframe=black!40,
    title=\textbf{Tier A --- Approaching Proficiency},
    fonttitle=\bfseries, arc=2mm, left=3mm, right=3mm, top=2mm, bottom=2mm, breakable]
\small
The backdrop's area is $A(x) = x^2+3x-28$ square feet and its width is $(x-4)$ feet.

\begin{enumerate}[leftmargin=1.6em, itemsep=4pt, topsep=4pt]
  \item Factor $A(x)$ completely.
\begin{work}
  A(x) &= x^2+3x-28 \\
       &= (x+7)(x-4)
\end{work}

  \item Find the backdrop's \textbf{length} by simplifying \ \TallMath{\frac{A(x)}{x-4}}. State the
        excluded value, and say why the shop teacher never has to worry about it.
\begin{work}
  \frac{A(x)}{x-4} &= \frac{(x+7)(x-4)}{x-4} = x+7 \quad \text{feet} \\
  x-4 &= 0 \; \Rightarrow \; x = 4 \; \text{excluded --- but a width must be positive}
\end{work}

  \item The shop teacher's drawing uses $x=9$. Give the backdrop's width, its length, and its area,
        then check your area against $A(9)$.
\begin{work}
  \text{width} &= 9-4 = 5 \text{ ft}, \qquad \text{length} = 9+7 = 16 \text{ ft} \\
  5 \cdot 16 &= 80 \text{ sq ft}, \qquad A(9) = 81+27-28 = 80 \quad\checkmark
\end{work}

  \item The paint order comes to $C(x) = 2x^2+6x-56$ dollars. Find the \textbf{cost per square
        foot} by simplifying \ \TallMath{\frac{C(x)}{A(x)}}, then say what that number means.
\begin{work}
  \frac{C(x)}{A(x)} &= \frac{2\left(x^2+3x-28\right)}{x^2+3x-28} = \frac{2(x+7)(x-4)}{(x+7)(x-4)} = 2 \\
  \text{at } x=9: \; \frac{160}{80} &= 2 \quad\checkmark
\end{work}
\ansline{Paint costs a flat \$2 per square foot --- the rate is the same for every backdrop size.}

  \item Multiply and simplify: \ \TallMath{\frac{x+3}{x-4} \cdot \frac{x^2-16}{x^2-9}}. State every
        excluded value.
\begin{work}
  \frac{x+3}{x-4} \cdot \frac{(x+4)(x-4)}{(x+3)(x-3)}
      &= \frac{x+4}{x-3} \qquad (x \ne 4, \; x \ne 3, \; x \ne -3)
\end{work}
\end{enumerate}
\end{tcolorbox}

\vspace{0.10in}

\boxguard[30]
\begin{tcolorbox}[colback=white, colframe=black!40, title=\textbf{Tier E --- Extension},
    fonttitle=\bfseries, arc=2mm, left=3mm, right=3mm, top=2mm, bottom=2mm, breakable]
\small

\begin{enumerate}[leftmargin=1.6em, itemsep=4pt, topsep=4pt]
  \item \textbf{Critique the work.} Two students turned in these claims. Neither is acceptable as
        written. Correct each one, then name the mistake.
        \[ \text{(a) } \frac{x+10}{x+5} = \frac{10}{5} = 2 \qquad\qquad
           \text{(b) } \frac{x^2-9}{x-3} = x+3 \text{ \emph{for every} } x \]
\begin{work}
  \text{(a) at } x=5: \; \frac{15}{10} &= \frac32 \ne 2 \; \text{--- } x \text{ is a \emph{term}, so nothing cancels} \\
  \text{(b) } \frac{(x+3)(x-3)}{x-3} &= x+3 \; \text{is right, but only for } x \ne 3 \\
  \text{at } x=3 \text{ the original reads } \tfrac00 \; &\text{(undefined) while } x+3 \text{ gives } 6
\end{work}
\ansline{(a) treats terms as factors; $\frac{x+10}{x+5}$ is already in simplest form.}\\
\ansline{(b) drops the scar: the simplification is correct only where the original was defined.}\\

  \item \textbf{Divide.} Simplify \ \TallMath{\frac{x^2-1}{2x+6} \div \frac{x+1}{x^2-9}} \ and state
        every excluded value --- including the ones that vanish along the way.
\begin{work}
  \frac{x^2-1}{2x+6} \cdot \frac{x^2-9}{x+1}
      &= \frac{(x+1)(x-1)}{2(x+3)} \cdot \frac{(x+3)(x-3)}{x+1} \\
      &= \frac{(x-1)(x-3)}{2} \qquad (x \ne -3, \; x \ne -1, \; x \ne 3)
\end{work}

  \item \textbf{Subtract, then justify.} Show that \ \TallMath{\frac{x^2}{x-5} - \frac{25}{x-5}} \
        simplifies to $x+5$. Then \emph{verify} your claim by evaluating both the original
        expression and $x+5$ at $x=7$, and state the one value of $x$ for which the two are
        \textbf{not} interchangeable.
\begin{work}
  \frac{x^2}{x-5} - \frac{25}{x-5} &= \frac{x^2-25}{x-5} = \frac{(x+5)(x-5)}{x-5} = x+5 \\
  \text{at } x=7: \; \frac{49}{2} - \frac{25}{2} &= \frac{24}{2} = 12, \qquad 7+5 = 12 \quad\checkmark
\end{work}
\ansline{At $x=5$: the original is undefined, while $x+5$ returns $10$. Everywhere else they agree.}
\end{enumerate}
\end{tcolorbox}

\end{document}
